% ============================================================================
%  APPENDIX: Annotation Codebook and LLM Annotation Prompts
%  Requires: \usepackage{listings} (for the prompt block) and booktabs.
%  If you cannot use listings, replace the lstlisting block with verbatim.
% ============================================================================
\section{Annotation Codebook and LLM Annotation Prompts}
\label{app:codebook}

This appendix documents the codebook shared by the human annotators and the local
LLM annotator, and the exact prompt used for the per-utterance annotation. Human
coders and the model applied the \emph{same} definitions (Tables~\ref{tab:taxonomy}--\ref{tab:levels}).

\subsection{Goal--Action Taxonomy}
\label{app:taxonomy}

Each reflective utterance is mapped to one or more item \emph{codes} rather than free
text. A code \texttt{G$g$.$i$} denotes a specific listed action; \texttt{G$g$.OTHER}
denotes an utterance that reflects goal $g$ but none of its listed actions; and
\texttt{OTHER} denotes a reflection that fits none of Goals~1--5. Codes map back to the
canonical action strings deterministically, so all outputs use the exact taxonomy wording.

\begin{table}[h]
\centering
\small
\caption{Goal--action taxonomy with item codes.}
\label{tab:taxonomy}
\begin{tabular}{@{}ll@{}}
\toprule
\textbf{Code} & \textbf{Action} \\
\midrule
\multicolumn{2}{@{}l}{\textit{Goal 1: Establish a Supportive and Professional Environment}} \\
G1.1 & Addressed family member by name. \\
G1.2 & Introduced him/herself by name and role. \\
G1.3 & Clearly stated the name of the deceased family member. \\
G1.4 & Sat down. (Body language / eye contact) \\
G1.5 & Displayed professional attire / presence. \\
G1.6 & Handled interruptions in non-disruptive manner. \\
G1.7 & Conducted interaction in organized manner. \\
\addlinespace
\multicolumn{2}{@{}l}{\textit{Goal 2: Assess and Align Expectations}} \\
G2.1 & Ensured all important survivors were present. \\
G2.2 & Determined knowledge survivors possessed. \\
G2.3 & Involved me when discussing reason for visit. \\
G2.4 & Elicited patient perspective of health situation. \\
\addlinespace
\multicolumn{2}{@{}l}{\textit{Goal 3: Deliver the News}} \\
G3.1 & Provided appropriate opening statement (warning shot). \\
G3.2 & Accurately / succinctly chronicled events leading to death. \\
G3.3 & Used phrase `dead' or `died' (avoided euphemisms). \\
G3.4 & Avoided jargon or explained terms. \\
\addlinespace
\multicolumn{2}{@{}l}{\textit{Goal 4: Manage the Emotional Response}} \\
G4.1 & Paused to allow family to assimilate information. \\
G4.2 & Responded to cues with appropriate touch. \\
G4.3 & Emotional response did not interfere with communication. \\
G4.4 & Legitimized my emotions. \\
G4.5 & Reinforced positive behaviors. \\
\addlinespace
\multicolumn{2}{@{}l}{\textit{Goal 5: Ensure Understanding and Facilitate Closure}} \\
G5.1 & Offered viewing of the deceased. \\
G5.2 & Established availability to answer questions. \\
G5.3 & Encouraged questions / concerns. \\
G5.4 & Summarized the interview. \\
G5.5 & Checked for accuracy during interview. \\
G5.6 & Reviewed next step(s). \\
G5.7 & Verified patient's understanding. \\
\addlinespace
\multicolumn{2}{@{}l}{\textit{Fallback codes}} \\
G$g$.OTHER & Reflects goal $g$ but none of its listed actions. \\
OTHER & Does not fit any of Goals 1--5. \\
\bottomrule
\end{tabular}
\end{table}

\subsection{Reflection Valence}
\label{app:valence}

Each reflection is assigned one valence, capturing how the speaker appraises the behavior.

\begin{table}[h]
\centering
\small
\caption{Valence codes.}
\label{tab:valence}
\begin{tabular}{@{}ll@{}}
\toprule
\textbf{Value} & \textbf{Definition} \\
\midrule
positive & Appraises the behavior as a strength / good choice. \\
negative & Appraises it as a shortcoming / mistake / regret / \\
         & something to change. \\
neutral  & Pure description, no positive or negative appraisal. \\
mixed    & Both a positive \emph{and} a negative in the same reflection. \\
\bottomrule
\end{tabular}
\end{table}

\subsection{Reflective Depth}
\label{app:levels}

Depth is coded on a five-level ordinal scale; coders assign the \emph{highest} level that applies.

\begin{table}[h]
\centering
\small
\caption{Reflective-depth levels (code the highest that applies).}
\label{tab:levels}
\begin{tabular}{@{}ll@{}}
\toprule
\textbf{Level} & \textbf{Definition} \\
\midrule
R0 Description & Only what happened, or a bare ``I did well'', with no reason. \\
R1 Reflective Description & What + why (gives a reason), but not questioned. \\
R2 Dialogic & Questions own reading / weighs an alternative / \\
            & reads family's hidden state. \\
R3 Transformative & Commits to change or a changed view of self \\
                  & (``next time I'll\ldots''). \\
R4 Critical & Beyond this encounter: culture / ethics / \\
            & training system (rare). \\
\bottomrule
\end{tabular}
\end{table}

\subsection{LLM Annotation Prompt}
\label{app:prompt}

The corpus was annotated per utterance with a local model (mistral-nemo:12b, Q4\_K\_M,
temperature~0). Each target utterance is judged within a $\pm2$-utterance sliding context
window, and the four decisions (\textit{is-reflection} $\rightarrow$ \textit{goal/action code}
$\rightarrow$ \textit{valence} $\rightarrow$ \textit{depth}) are made in a single prompt that
returns strict JSON. Pure backchannels / filler ($\leq 3$ words) are filtered before the model
call. The template is shown below; \texttt{\{...\}} placeholders are filled at run time (the
taxonomy from Table~\ref{tab:taxonomy}, valence from Table~\ref{tab:valence}, and depth from
Table~\ref{tab:levels} are injected verbatim).

\begin{lstlisting}[basicstyle=\ttfamily\footnotesize,breaklines=true,frame=single]
You are an expert qualitative researcher analyzing a medical debriefing
transcript after a BREAKING BAD NEWS simulation. Speakers are now
reflecting on how it went.

TAXONOMY (map to an item CODE below; do NOT invent or paraphrase item text):
{goal-action taxonomy, Table 1}

CONVERSATION CONTEXT:
[2 ago] {role}: "{utterance}"
[1 ago] {role}: "{utterance}"
---
[TARGET UTTERANCE] {role}: "{utterance}"
---
[next] {role}: "{utterance}"

RULES:
1. ONLY REFLECTIONS: annotate the TARGET only when the speaker
   evaluates/recalls/reflects on a clinical action from the simulation.
   Otherwise set "is_reflection" false and return empty lists.
2. IGNORE SMALL TALK: filler ("yeah", "okay"), logistics, bare
   acknowledgements -> is_reflection false.
3. INHERIT CONTEXT: a generic agreement ("exactly") right after a
   specific reflection inherits it.
4. MAP TO ITEM CODES: first reason which GOAL the utterance really shows
   (e.g. silence/pauses and touch are emotional-response actions, NOT
   setup or assessment), THEN put the taxonomy code(s) in
   "reflected_item_codes".
   - Each entry MUST be one of the codes in the taxonomy above (codes
     span G1.1 through G5.x, plus G#.OTHER and OTHER). Do NOT default to
     a goal's first item; pick the item that matches.
   - Do NOT write item text and do NOT invent new codes.
   - CONSISTENCY: each code must be supported by the words shown in the
     CONTEXT WINDOW (the target plus the surrounding utterances). If
     nothing in that window clearly expresses the action, pick a
     different code or set is_reflection false.
   - Usually one code; use several only if the utterance clearly
     reflects on several actions.
5. When an OTHER code (G#.OTHER or OTHER) is used, write a 3-7 word
   "other_summary".
6. "valence" = ONE of ["positive","negative","neutral","mixed"] (or
   "None" if not a reflection):
{valence definitions, Table 2}
7. "level" = ONE of ["R0","R1","R2","R3","R4"] (CODE THE HIGHEST THAT
   APPLIES; "None" if not a reflection):
{depth definitions, Table 3}

OUTPUT (Strict JSON):
{
    "is_reflection": true or false,
    "reason": "which goal the utterance shows and why (decide this
               BEFORE the codes)",
    "reflected_item_codes": ["<best code(s) from the taxonomy above>"],
    "valence": "positive|negative|neutral|mixed|None",
    "level": "R0|R1|R2|R3|R4|None",
    "other_summary": "short summary if Other else empty"
}
\end{lstlisting}

The returned codes are mapped back to canonical (goal, action) pairs deterministically;
an utterance that is not a reflection, or whose code list is empty, is recorded as
non-reflective with null valence and depth.
